%% sections/04-evaluation.tex

\section{Cross-Campaign Validation}
\label{sec:evaluation}

The central empirical claim of the player-style framework is that ratified styles are
\textit{not archetype-specific}: they fire across campaign archetypes with different
motivational structures, party compositions, and designed emotional registers. To test
this claim, we formulated 12 research validation hypotheses (RV1--RV12) and evaluated
each against the 7-campaign corpus.

\subsection{Research Hypotheses}

The 12 hypotheses are organized into three groups.

\paragraph{Group 1: Professional trigger campaigns (RV1--RV4; Campaign 5).}
Campaign 5 was designed as a Thieves Guild campaign: professional motivations, no
designed grief hooks, PCs with institutional loyalties rather than personal losses.
The group tests whether styles that originated in grief-themed campaigns (C1, C3)
transfer to a professional-motivation context.

\begin{description}
  \item[RV1.] Sheet-deep-reader fires in a campaign with no designed grief hooks.
  \item[RV2.] Craft-witness fires from professional practice as reliably as from
    religious or craft-ritual practice.
  \item[RV3.] NPC-arc-completion trigger fires in a campaign where no NPC arc was
    designed around grief.
  \item[RV4.] At least one PC-authored finale deliverable is constructed and deployed
    in a campaign with professional motivations.
\end{description}

\paragraph{Group 2: Duty trigger campaigns (RV5--RV8; Campaign 6).}
Campaign 6 was designed as a Military Unit campaign: duty-based institutional
motivations, hierarchical structure, a party defined by rank and obligation rather
than personal choice. The group tests styles in the absence of personal-stakes design.

\begin{description}
  \item[RV5.] Sheet-deep-reader fires in a military-unit campaign (institutional
    training backgrounds, orders-received hooks).
  \item[RV6.] Act-without-announcement fires in a campaign where naming events
    are duty-related rather than grief-related.
  \item[RV7.] NPC-arc-completion trigger fires in a campaign where NPC arcs are
    structured around duty and dereliction rather than loss.
  \item[RV8.] Verbal-documentation-arc fires in a campaign where the PC's documentation
    practice is formal reporting rather than oral testimony.
\end{description}

\paragraph{Group 3: Situational-to-relationship trigger shift (RV9--RV12; Campaign 7).}
Campaign 7 was designed as a Party of Strangers campaign: no designed personal stakes,
no grief paragraphs, no professional codes, PCs with minimal backstory. The group
tests styles at maximum design sparsity, and tests the hypothesis that trigger type
shifts from situational (responding to external events) to relationship-driven
(responding to other PCs) across the campaign's arc.

\begin{description}
  \item[RV9.] Player styles fire without any designed personal stakes, provided
    character sheets are designed with sufficient specificity in other dimensions.
  \item[RV10.] Craft-witness fires from regional or professional knowledge habits
    rather than craft-ritual practice when craft-ritual is not designed in.
  \item[RV11.] Trigger type shifts from situational (Sessions 1--3) to
    relationship-driven (Sessions 4--7) as the strangers develop in-campaign history.
  \item[RV12.] Sheet-deep-reader instances in late sessions draw on in-campaign
    history rather than pre-campaign backstory, as the campaign itself becomes the
    ``sheet'' the PC reads.
\end{description}

\subsection{Results: Group 1 (Campaign 5, Professional Triggers)}

\textbf{RV1: Confirmed.} Sheet-deep-reader produced 5 instances in Campaign 5, all
drawing on institutional-training content rather than grief-paragraph content. Mira
Calloway's professional-code heuristic (``log what you observe; surface it only when
the framework is complete'') and Pale's organizational-chart knowledge produced
3 instances each. The style fired from professional sheet content in the absence of
grief hooks.

\textbf{RV2: Confirmed.} Craft-witness produced 3 instances in Campaign 5 from
Dace Morrow's locksmith-practice habit (assessing lock-wear patterns as a daily
professional attention). The practice is secular and professional; the instances are
structurally identical to religious-practice instances from Campaign 1. The practice
type is irrelevant; specificity is what matters.

\textbf{RV3: Confirmed.} NPC-arc-completion trigger produced 5 instances in Campaign 5.
Two NPC arcs were structured around professional betrayal (Harren's informant-arc)
and institutional exit (Cassel Vane's extraction-arc). Both fired by character logic
when party behavior met the specified firing conditions.

\textbf{RV4: Confirmed.} Mira Calloway's false dossier (MARTEN-7) was constructed
across 7 sessions and deployed at the finale. The deliverable is professional (a
fabricated intelligence document) rather than confessional or narrative, confirming
that the PC-authored-finale-deliverable pattern is not limited to emotionally resonant
content types.

\subsection{Results: Group 2 (Campaign 6, Duty Triggers)}

\textbf{RV5: Confirmed.} Sheet-deep-reader produced 3 instances in Campaign 6, all
drawing on military training and orders-received content. Sergeant Vann's tactical
assessment (``prior unit tactics inform terrain reading: three-exit approach'')
produced 2 instances; Lira Coldmarch's supply-accountability heuristic produced 1.
The institutional-training dimension of the style fires reliably even when the sheet
contains no personal backstory.

\textbf{RV6: Confirmed.} Act-without-announcement produced 2 instances in Campaign 6,
both duty-related. In Session 4, a PC named a dereliction of duty to a superior officer
(a naming event). In Session 5, the same PC moved to the rear of the formation when
entering the relevant type of terrain, without explanation, without announcement.
The behavior followed the naming; the PC did not declare the change.

\textbf{RV7: Confirmed.} NPC-arc-completion trigger produced 6 instances in Campaign 6.
Five NPC arcs were structured around duty, dereliction, or institutional loyalty; all
five fired by character logic when party behavior met conditions. Campaign 6 produced
the highest NPC-arc-completion instance count per session of any campaign in the corpus.

\textbf{RV8: Partially confirmed.} Verbal-documentation-arc produced 1 candidate
instance in Campaign 6 (a formal-reporting PC whose vocabulary broke down at a tribunal
scene). The single instance does not cross the confirmation threshold for Campaign 6
specifically, but it is consistent with the style's design trigger. The hypothesis is
marked partially confirmed pending further evidence.

\subsection{Results: Group 3 (Campaign 7, Situational-to-Relationship Shift)}

\textbf{RV9: Confirmed.} Four styles fired in Campaign 7 without any designed
personal stakes. Sheet-deep-reader fired from regional knowledge (Brek Ashfall's
terrain-advantage instances), craft-witness fired from structural habit (Vess
Ironmark's load-path assessment), act-without-announcement fired from a relational
naming event (Tam's road-companion behavior following Brek's wrong-fork decision),
and npc-arc-completion trigger fired in 5 instances. The styles do not require grief
hooks; they require design specificity, wherever that specificity is located.

\textbf{RV10: Confirmed.} Craft-witness produced 2 instances in Campaign 7 from
Vess Ironmark's structural-assessment habit (``door set wrong; frame doesn't carry
the load this passage suggests''). This is regional-professional knowledge, not
craft-ritual. The instances are structurally identical to religious-practice instances.

\textbf{RV11: Confirmed.} A trigger-type analysis across Campaign 7's sessions
shows a measurable shift in what causes style instances to fire (Table~\ref{tab:trigger-shift}).
Sessions 1--3 produced primarily situational triggers (a scene presents an opportunity
the PC's sheet is configured to notice). Sessions 4--7 produced primarily relationship-driven
triggers (a prior in-campaign interaction with another PC becomes the context for the
instance). The shift is not absolute --- situational triggers continue in Sessions 4--7
--- but the proportion of relationship-driven triggers rises from 14\% in Sessions 1--3
to 61\% in Sessions 4--7.

\textbf{RV12: Confirmed.} In Campaign 7 Sessions 5--7, the majority of
sheet-deep-reader instances drew on in-campaign history (things that had happened in
prior sessions) rather than pre-campaign backstory. In Sessions 1--3, all
sheet-deep-reader instances drew on pre-campaign content (the sparse backstory on the
sheets). By Session 6, Brek Ashfall's sheet-deep-reader instances were drawing
primarily on decisions made in Sessions 2--4 as the relevant ``sheet'' content.
The campaign itself becomes the character sheet.

\subsection{Aggregate Instance Counts}

Table~\ref{tab:cross-campaign-counts} presents per-style, per-campaign confirmed
instance counts across all 7 campaigns.

\begin{table*}[t]
\centering
\caption{Cross-Campaign Confirmed Instance Counts by Style}
\label{tab:cross-campaign-counts}
\begin{tabular}{@{}lrrrrrrrr@{}}
\toprule
Style & C1 & C2 & C3 & C4 & C5 & C6 & C7 & Total \\
\midrule
sheet-deep-reader           & 13 & 12 &  8 &  9 &  5 &  3 &  4 &  54 \\
craft-witness               & 16 & 14 &  6 &  4 &  3 &  2 &  2 &  47 \\
npc-arc-completion trigger  &  7 & 14 &  8 &  7 &  5 &  6 &  5 &  52 \\
act-without-announcement    &  1 & 11 &  4 &  2 &  0 &  0 &  1 &  19 \\
verbal-documentation-arc    &  0 &  5 &  7 &  2 &  0 &  1 &  0 &  15 \\
documentary-witness         &  0 &  0 &  0 &  7 &  0 &  0 &  4 &  11 \\
act-without-ann.\ intra-session &  0 &  3 &  2 &  2 &  0 &  0 &  1 &   8 \\
\midrule
\textbf{Total per campaign} & \textbf{37} & \textbf{59} & \textbf{35} & \textbf{33} & \textbf{13} & \textbf{12} & \textbf{17} & \textbf{206} \\
\bottomrule
\end{tabular}
\end{table*}

\begin{table}[t]
\centering
\caption{Trigger-Type Breakdown for Campaign 7: Sessions 1--3 vs.\ Sessions 4--7}
\label{tab:trigger-shift}
\begin{tabular}{@{}lrrrr@{}}
\toprule
 & \multicolumn{2}{c}{Sessions 1--3} & \multicolumn{2}{c}{Sessions 4--7} \\
\cmidrule(lr){2-3} \cmidrule(lr){4-5}
Trigger type & Count & \% & Count & \% \\
\midrule
Situational (scene-driven)      & 6 & 86\% &  4 & 39\% \\
Relationship-driven (PC-to-PC)  & 1 & 14\% &  6 & 61\% \\
\midrule
\textbf{Total}                  & 7 & 100\% & 10 & 100\% \\
\bottomrule
\end{tabular}
\end{table}

\subsection{Research Hypothesis Summary}

Table~\ref{tab:rv-summary} summarizes the outcome for all 12 hypotheses.

\begin{table}[t]
\centering
\caption{Research Validation Hypotheses: Outcomes}
\label{tab:rv-summary}
\begin{tabular}{@{}llp{4cm}l@{}}
\toprule
Hypothesis & Group & Claim (abbreviated) & Result \\
\midrule
RV1  & C5 professional & Sheet-deep-reader without grief hooks & Confirmed \\
RV2  & C5 professional & Craft-witness from professional practice & Confirmed \\
RV3  & C5 professional & NPC-arc-completion without grief-arc design & Confirmed \\
RV4  & C5 professional & PC-authored deliverable (professional) & Confirmed \\
RV5  & C6 duty         & Sheet-deep-reader from institutional training & Confirmed \\
RV6  & C6 duty         & Act-without-announcement (duty naming) & Confirmed \\
RV7  & C6 duty         & NPC-arc-completion (duty arcs) & Confirmed \\
RV8  & C6 duty         & Verbal-documentation-arc (formal reporting) & Partial \\
RV9  & C7 situational  & Styles without personal stakes & Confirmed \\
RV10 & C7 situational  & Craft-witness from regional knowledge & Confirmed \\
RV11 & C7 situational  & Trigger shift: situational to relationship & Confirmed \\
RV12 & C7 situational  & Sheet-deep-reader from in-campaign history & Confirmed \\
\midrule
\multicolumn{3}{l}{\textbf{Confirmed / Partial / Disconfirmed}} & 11 / 1 / 0 \\
\bottomrule
\end{tabular}
\end{table}

Eleven of 12 hypotheses were fully confirmed; RV8 was partially confirmed (1 candidate
instance, insufficient for Campaign 6 confirmation). No hypothesis was disconfirmed.
The partial confirmation of RV8 is a finding in itself: the verbal-documentation-arc
style may require oral-practice design in the PC sheet to fire reliably, and the formal-reporting
variant may be a weaker trigger than the oral-testimony variant.
